\chapter{Classifying the reduced-boundary blow-up}\label{sec:perimeter-halfspace-classification}
Throughout this section $n\geq2$, and $E$ has locally finite perimeter in an open set $U\subset\mathbb R^n$.
We now classify every subsequential limit using distributional derivatives and the two-phase volume bound. The argument follows the blow-up part of \cite{Sim18}, Chapter 3, Theorem 4.4, printed pp.~98--99, with all signs derived from the project's outward convention $D\chi_E=-\nu\,|D\chi_E|$. The limiting set will be $\{y:y\cdot\nu_E(x)<0\}$.

\begin{lemma}[Zero tangential derivatives imply one-dimensional dependence]\label{lem:distributional-one-direction}
Suppose $u\in L^\infty(\mathbb R^n)$ and $D_i u=0$ in distributions for $i=1,\ldots,n-1$. Then there is $h\in L^\infty(\mathbb R)$ such that
\[
 u(x',s)=h(s)\quad\text{for almost every }(x',s)\in\mathbb R^{n-1}\times\mathbb R.
\]
If $u$ is a characteristic function, $h$ is zero or one almost everywhere. This makes explicit the cylindrical-limit step after equation (8) in \cite{Sim18}, p.~98.
\end{lemma}
\begin{proof}
Convolve $u$ with a standard smooth compactly supported approximate identity on $\mathbb R^n$. Boundedness makes the convolution well defined even when $u$ is not globally integrable. Each smooth $u_\delta$ has zero derivative in every $x'$ direction, hence is constant on each connected horizontal plane: $u_\delta(x',s)=h_\delta(s)$. On any bounded product $Q\times I$, with $0<|Q|<\infty$, local $L^1$ convergence of convolutions implies
\[
 |Q|\int_I|h_\delta-h_\gamma|\,ds
 =\int_{Q\times I}|u_\delta-u_\gamma|\,dx\longrightarrow0.
\]
Thus $h_\delta$ has an $L^1_{\mathrm{loc}}(\mathbb R)$ limit $h$. Limits on overlapping intervals agree, and the uniform bound passes to $h$ by almost-everywhere convergence along a sequence. On every product cylinder $u=h(s)$ in $L^1$, hence almost everywhere; a countable exhaustion gives the global assertion. If $u$ is zero or one almost everywhere, Fubini on any such product shows that $h$ has the same property.
\end{proof}

\begin{lemma}[A one-dimensional monotone phase is a half-space]\label{lem:monotone-phase-halfspace}
Suppose $F\subset\mathbb R^n$ is measurable, $D_i\chi_F=0$ for $i<n$, and $D_n\chi_F\leq0$ in the distributional sense. Then for some $\lambda\in[-\infty,\infty]$,
\[
 F=\{(x',s):s<\lambda\}\quad\text{up to a volume-null set}.
\]
The endpoint values mean the empty set and all of $\mathbb R^n$. This classification requires neither a smooth boundary nor a prior rectifiability theorem.
\end{lemma}
\begin{proof}
By Lemma \ref{lem:distributional-one-direction}, $\chi_F=h(s)$ with $h\in\{0,1\}$ almost everywhere. Test the derivative sign with $\phi(x',s)=a(x')b(s)$, where $a\geq0$ is smooth, compactly supported and has integral one, and $b\geq0$ is a smooth scalar test. It gives $Dh\leq0$ in distributions. Nonnegative one-dimensional mollification therefore gives smooth nonincreasing functions $h_\delta$, bounded between zero and one, converging to $h$ locally in $L^1$.

Choose a sequence of mollification radii with almost-everywhere convergence on every bounded interval. On the full-measure set of convergence points, the inequality $h(s)\geq h(t)$ holds whenever $s<t$, by passing to the limit in the corresponding smooth inequalities. Remove also the null set where $h$ is not zero or one. The set of remaining points with value one is an initial segment of this ordered full-measure set, and those with value zero form its final segment. Their separating endpoint is $\lambda$, allowing infinite endpoints. Consequently $h=\chi_{(-\infty,\lambda)}$ almost everywhere, which proves the claim. This argument avoids assigning an arbitrary pointwise representative at the jump.
\end{proof}

\begin{proposition}[Compactness of every sequence of reduced-boundary dilations]\label{prop:reduced-boundary-blowup-compactness}
Let $x\in\partial^*E$ and let $r_j\downarrow0$ be any sequence. A subsequence of $\chi_{E_{x,r_j}}$ converges in $L^1_{\mathrm{loc}}(\mathbb R^n)$ to $\chi_F$ for a set $F$ of locally finite perimeter. Its rescaled perimeter measures satisfy, for each fixed $R>0$ and all sufficiently large $j$,
\[
 |D\chi_{E_{x,r_j}}|(B_R)\leq2n\omega_n R^{n-1}.
\]
This is the compactness passage preceding equation (7) in \cite{Sim18}, p.~98.
\end{proposition}
\begin{proof}
By Lemma \ref{lem:perimeter-blowup-scaling} the left side is $r_j^{1-n}\mu(B_{r_jR}(x))$. Proposition \ref{prop:reduced-boundary-perimeter-growth} gives the displayed bound as soon as $r_jR<r_0$. The characteristic functions have $L^1$ norm at most $|B_R|$ on that ball. Use Corollary \ref{cor:perimeter-compactness} on successively larger balls and extract a diagonal subsequence. The dilated domains contain each fixed ball eventually; discard a finite initial segment before applying compactness on that ball. This does not require finite perimeter outside the original open set. The local limits agree on overlaps and are characteristic functions. Lower semicontinuity gives local finite perimeter of $F$.
\end{proof}

\begin{proposition}[Polar alignment determines the direction and sign of a limit]\label{prop:reduced-boundary-limit-direction}
Every limit $F$ in Proposition \ref{prop:reduced-boundary-blowup-compactness} agrees almost everywhere with
\[
 \{y:y\cdot\nu_E(x)<\lambda\}
\]
for some $\lambda\in[-\infty,\infty]$. The derivative of its characteristic function in the outward direction is nonpositive. This is the sign-sensitive part of the classification on printed p.~98 of \cite{Sim18}.
\end{proposition}
\begin{proof}
Translate $x$ to zero and rotate so that $\nu_0=e_n$. Put $\mu_j=|D\chi_{E_{0,r_j}}|$ and $\nu_j(y)=\nu(r_jy)$. On each fixed ball $B_R$, Corollary \ref{cor:perimeter-polar-components} and the scaling lemma give
\[
 \int_{B_R}|P\nu_j|\,d\mu_j
 \leq\sqrt{\varepsilon(r_jR)}\,\mu_j(B_R)\longrightarrow0,
 \qquad
 \int_{B_R}(\nu_j\cdot e_n)_-\,d\mu_j\longrightarrow0.
\]
The bounded masses from Proposition \ref{prop:reduced-boundary-blowup-compactness} justify both limits. By Proposition \ref{prop:bv-derivative-convergence}, $D\chi_{E_{0,r_j}}\rightharpoonup D\chi_F$ locally. Since $D\chi_{E_{0,r_j}}=-\nu_j\mu_j$, the first estimate makes every horizontal component of $D\chi_F$ zero. For a nonnegative compactly supported scalar test $\phi$ in $B_R$,
\[
 \int\phi\,dD_n\chi_F
 =\lim_j-\int\phi\,\nu_j\cdot e_n\,d\mu_j
 \leq\limsup_j\|\phi\|_\infty
                    \int_{B_R}(\nu_j\cdot e_n)_-\,d\mu_j=0.
\]
This is a sign assertion about the limit; it does not assume each weighted integral already has the correct sign at finite scale. Lemma \ref{lem:monotone-phase-halfspace} gives the asserted half-space. Rotate and translate back to the stated rescaled coordinates.
\end{proof}

\begin{proposition}[The two-phase bounds center the limiting half-space]\label{prop:reduced-boundary-centered-limit}
The number $\lambda$ in Proposition \ref{prop:reduced-boundary-limit-direction} is zero. Thus every subsequential limit is
\[
 H_x=\{y:y\cdot\nu_E(x)<0\}.
\]
This is the role of equations (10)--(11) in \cite{Sim18}, p.~99.
\end{proposition}
\begin{proof}
For any fixed $R>0$, scaling Theorem \ref{thm:reduced-boundary-two-phase-volume} gives
\[
 |E_{x,r_j}\cap B_R|\geq c_nR^n,
 \qquad |B_R\setminus E_{x,r_j}|\geq c_nR^n
\]
for all sufficiently large $j$. Local $L^1$ convergence passes these two inequalities to $F$ on $B_R$. They hold for every fixed $R>0$. The empty and full sets violate one of them. A half-space with finite $\lambda\ne0$ is empty or full inside $B_R$ whenever $0<R<|\lambda|$, again a contradiction. Therefore $\lambda=0$.
\end{proof}

\begin{theorem}[Full local convergence to the oriented half-space]\label{thm:reduced-boundary-full-blowup}
For every $x\in\partial^*E$, as $r\downarrow0$,
\[
 \chi_{E_{x,r}}\longrightarrow\chi_{H_x}
 \quad\text{in }L^1_{\mathrm{loc}}(\mathbb R^n),
 \qquad H_x=\{y:y\cdot\nu_E(x)<0\}.
\]
This is the characteristic-function conclusion of \cite{Sim18}, Chapter 3, Theorem 4.4, proved without invoking the rectifiability conclusion of that theorem.
\end{theorem}
\begin{proof}
Suppose convergence failed on some ball $B_R$. There would be a positive $\delta$ and a sequence $r_j\downarrow0$ for which $\int_{B_R}|\chi_{E_{x,r_j}}-\chi_{H_x}|\geq\delta$. Proposition \ref{prop:reduced-boundary-blowup-compactness} supplies a locally convergent subsequence. Propositions \ref{prop:reduced-boundary-limit-direction} and \ref{prop:reduced-boundary-centered-limit} identify its limit with $H_x$, contradicting the lower bound. Since this works for every ball, it proves full local convergence and not merely convergence of one selected sequence.
\end{proof}
